\documentclass[aspectratio=169,10pt]{beamer}

% -------------------------------------------------
% Pacotes basicos
% -------------------------------------------------
\usepackage[utf8]{inputenc}
\usepackage[T1]{fontenc}
\usepackage[brazil]{babel}
\usepackage{lmodern}
\usepackage{amsmath,amssymb}
\usepackage{booktabs}
\usepackage{array}
\usepackage{tikz}
\usepackage{graphicx}
\usetikzlibrary{arrows.meta,positioning,shapes.geometric,fit,calc,backgrounds}
\usepackage{multicol}
\usepackage{ragged2e}

% -------------------------------------------------
% Estilo visual inspirado no modelo enviado
% -------------------------------------------------
\definecolor{InovGreen}{RGB}{68,40,155}
\definecolor{InovLightGreen}{RGB}{242,239,253}
\definecolor{InovDark}{RGB}{40,36,55}
\definecolor{InovGray}{RGB}{248,247,252}
\definecolor{InovOrange}{RGB}{122,113,245}
\definecolor{InovBlue}{RGB}{92,116,255}
\definecolor{InovPurple}{RGB}{68,40,155}
\definecolor{InovLilac}{RGB}{242,239,253}

\setbeamertemplate{navigation symbols}{}
\setbeamertemplate{itemize item}{\color{InovGreen}$\blacktriangleright$}
\setbeamertemplate{itemize subitem}{\color{InovOrange}$\bullet$}
\setbeamertemplate{enumerate item}{\color{InovGreen}\insertenumlabel.}
\setbeamercolor{title}{fg=InovGreen}
\setbeamercolor{frametitle}{fg=InovGreen,bg=white}
\setbeamercolor{block title}{fg=white,bg=InovGreen}
\setbeamercolor{block body}{fg=InovDark,bg=InovLightGreen}
\setbeamercolor{alerted text}{fg=InovOrange}

\setbeamertemplate{frametitle}{%
  \vspace{-0.05cm}
  \insertframetitle
  \par\vspace{-0.1cm}
  {\color{InovGreen}\rule{\linewidth}{0.5pt}}
}

\setbeamertemplate{footline}{%
  \begin{tikzpicture}[remember picture,overlay]
    \node[anchor=south east, xshift=-0.35cm, yshift=0.18cm] at (current page.south east)
      {\color{InovDark}\scriptsize \insertframenumber/\inserttotalframenumber};
  \end{tikzpicture}%
}

% -------------------------------------------------
% Macros para diagramas
% -------------------------------------------------
\tikzset{
  box/.style={rectangle, rounded corners=3pt, draw=InovGreen, very thick, align=center, fill=InovLightGreen, minimum height=0.75cm, text width=2.6cm},
  smallbox/.style={rectangle, rounded corners=3pt, draw=InovGreen, thick, align=center, fill=InovGray, minimum height=0.65cm, text width=2.2cm},
  orangebox/.style={rectangle, rounded corners=3pt, draw=InovOrange, very thick, align=center, fill=orange!12, minimum height=0.75cm, text width=2.6cm},
  bluebox/.style={rectangle, rounded corners=3pt, draw=InovBlue, very thick, align=center, fill=blue!8, minimum height=0.75cm, text width=2.6cm},
  arrow/.style={-{Latex[length=2.5mm]}, very thick, draw=InovDark},
  dashedarrow/.style={-{Latex[length=2.5mm]}, thick, dashed, draw=InovOrange}
}

\title[Pipeline PepMem-AI]{Pipeline de IA para Predição de Interação Peptídeo--Membrana}
\subtitle{Bioprospecção de peptídeos escorpiônicos com embeddings, modelos multimodais e validação experimental}
\author{Prof. Marcelo A. C. Fernandes \; | \; InovAI Lab}
\institute{UFRN \; | \; IMD \; | \; LANCE/InovAI Lab}
\date{2026}

% Fundo do InovAI Lab utilizado em todos os slides
\newcommand{\inovaibackground}{%
  \begin{tikzpicture}[remember picture,overlay]
    \node[anchor=south west, inner sep=0pt] at (current page.south west) {%
      \includegraphics[width=\paperwidth,height=\paperheight]{/mnt/data/inovai_lab_background.png}%
    };
    % logo pequeno no canto superior direito com fundo branco semitransparente
    \node[anchor=north east, xshift=-0.16cm, yshift=-0.14cm,
      rounded corners=2pt, fill=white, fill opacity=0.72, text opacity=1,
      inner sep=1.5pt] at (current page.north east) {%
      \includegraphics[width=0.78cm]{/mnt/data/tech_meets_biology_logo_design.png}%
    };
  \end{tikzpicture}%
}
\usebackgroundtemplate{\inovaibackground}

% Variante do fundo sem logo no canto (para o slide de abertura)
\newcommand{\inovaibackgroundnocorner}{%
  \begin{tikzpicture}[remember picture,overlay]
    \node[anchor=south west, inner sep=0pt] at (current page.south west) {%
      \includegraphics[width=\paperwidth,height=\paperheight]{/mnt/data/inovai_lab_background.png}%
    };
  \end{tikzpicture}%
}

\begin{document}

% -------------------------------------------------
{\usebackgroundtemplate{\inovaibackgroundnocorner}
\begin{frame}[plain]
  \vspace{0.15cm}
  \begin{center}
  \begin{tikzpicture}
    \node[rounded corners=10pt, fill=white, fill opacity=0.90, text opacity=1,
          draw=InovPurple, line width=0.9pt, inner sep=0.40cm,
          text width=0.80\paperwidth, align=center] {
      {\Large\bfseries\color{InovPurple} Pipeline de IA para Predição de Interação Peptídeo--Membrana}\\[0.22cm]
      {\large\color{InovDark} Bioprospecção de peptídeos escorpiônicos com validação experimental}\\[0.18cm]
      \includegraphics[width=0.17\paperwidth]{/mnt/data/tech_meets_biology_logo_design.png}\\[0.12cm]
      {\normalsize Prof. Marcelo A. C. Fernandes}\\[0.05cm]
      {\normalsize InovAI Lab \; | \; IMD/UFRN}
    };
  \end{tikzpicture}
  \end{center}
\end{frame}
}

% -------------------------------------------------
\begin{frame}{Contexto do projeto}
\begin{block}{Projeto de bioprospecção}
O projeto busca identificar peptídeos escorpiônicos com potencial terapêutico, combinando triagem computacional, inteligência artificial e validação experimental in vitro.
\end{block}

\vspace{0.15cm}
\begin{columns}[T]
\column{0.52\textwidth}
\textbf{Candidatos principais}
\begin{itemize}
  \item Peptídeos derivados de venenos de escorpiões.
  \item Análogos de Stigmurina e TsAP-2.
  \item Moléculas com carga, hidrofobicidade e anfipaticidade favoráveis.
\end{itemize}

\column{0.48\textwidth}
\textbf{Atividades de interesse}
\begin{itemize}
  \item Antibacteriana e antifúngica.
  \item Antiparasitária.
  \item Antiviral.
  \item Antitumoral e baixa citotoxicidade.
\end{itemize}
\end{columns}
\end{frame}

% -------------------------------------------------
\begin{frame}{Por que usar IA no projeto?}
\begin{block}{Desafio experimental}
Testar todos os peptídeos contra todos os alvos biológicos pode ser caro, lento e pouco escalável.
\end{block}

\vspace{0.25cm}
\begin{center}
\resizebox{0.90\textwidth}{!}{%
\begin{tikzpicture}[node distance=0.8cm]
  \node[box] (pep) {Peptídeos\ candidatos};
  \node[box,right=of pep] (targets) {Múltiplas\ membranas-alvo};
  \node[orangebox,right=of targets] (ai) {Triagem\ por IA};
  \node[bluebox,right=of ai] (rank) {Ranking\ para bancada};
  \node[box,right=of rank] (lab) {Validação\ in vitro};
  \draw[arrow] (pep) -- (targets);
  \draw[arrow] (targets) -- (ai);
  \draw[arrow] (ai) -- (rank);
  \draw[arrow] (rank) -- (lab);
\end{tikzpicture}}
\end{center}

\vspace{0.2cm}
\begin{itemize}
  \item A IA prioriza os candidatos com maior chance de atividade e seletividade.
  \item A bancada confirma as predições e gera novos dados para melhorar o modelo.
\end{itemize}
\end{frame}

% -------------------------------------------------
\begin{frame}{Como chegar ao problema peptídeo--membrana}
\begin{center}
\resizebox{0.95\textwidth}{!}{%
\begin{tikzpicture}[node distance=0.8cm]
  \node[box,text width=2.8cm] (bio) {Peptídeos\ bioativos};
  \node[box,text width=3.0cm,right=of bio] (prop) {Propriedades\ físico-químicas};
  \node[orangebox,text width=3.0cm,right=of prop] (mem) {Interação com\ membranas};
  \node[bluebox,text width=3.0cm,right=of mem] (effect) {Efeito\ biológico};
  \draw[arrow] (bio) -- (prop);
  \draw[arrow] (prop) -- (mem);
  \draw[arrow] (mem) -- (effect);
\end{tikzpicture}}
\end{center}

\vspace{0.25cm}
\begin{block}{Reinterpretação computacional}
Do ponto de vista de IA, a pergunta central não é apenas se o peptídeo é ativo. A pergunta é: \textbf{como um peptídeo interage com uma membrana específica e qual efeito essa interação produz?}
\end{block}

\vspace{0.15cm}
\begin{itemize}
  \item Membranas bacterianas, fúngicas, parasitárias, virais, normais e tumorais são diferentes.
  \item O mesmo peptídeo pode ser ativo em uma membrana e tóxico ou inativo em outra.
\end{itemize}
\end{frame}

% -------------------------------------------------
\begin{frame}{Objetivo}
\begin{block}{Pergunta central}
Construir um pipeline computacional para predizer se um peptídeo escorpiônico interage, permeabiliza ou desestabiliza uma membrana biológica específica.
\end{block}

\vspace{0.2cm}
\begin{itemize}
  \item Formular o problema como \alert{interação peptídeo--membrana}.
  \item Representar peptídeos por descritores físico-químicos e embeddings.
  \item Representar membranas-alvo por carga, composição, estrutura e propriedades físicas.
  \item Construir modelos de IA para predição, ranqueamento e seleção experimental.
  \item Fechar o ciclo com validação in vitro e realimentação do modelo.
\end{itemize}
\end{frame}

% -------------------------------------------------
\begin{frame}{Motivação biológica}
\begin{columns}[T]
\column{0.53\textwidth}
\begin{itemize}
  \item Peptídeos de venenos de escorpiões podem apresentar atividades antimicrobiana, antifúngica, antiviral, antiparasitária e antitumoral.
  \item Muitas dessas atividades dependem da interação com membranas biológicas.
  \item Características importantes dos peptídeos:
    \begin{itemize}
      \item carga positiva;
      \item hidrofobicidade;
      \item anfipaticidade;
      \item tendência helicoidal;
      \item momento hidrofóbico.
    \end{itemize}
\end{itemize}

\column{0.47\textwidth}
\begin{center}
\begin{tikzpicture}[node distance=0.5cm]
  \node[box,text width=2.4cm] (p) {Peptídeo\ catiônico};
  \node[box,text width=2.4cm,below=of p] (m) {Membrana\ negativa};
  \node[orangebox,text width=2.8cm,below=of m] (i) {Interação\ eletrostática};
  \node[bluebox,text width=2.8cm,below=of i] (e) {Inserção, poro\ ou ruptura};
  \draw[arrow] (p) -- (m);
  \draw[arrow] (m) -- (i);
  \draw[arrow] (i) -- (e);
\end{tikzpicture}
\end{center}
\end{columns}
\end{frame}

% -------------------------------------------------
\begin{frame}{Formulação do problema}
\vspace{-0.35cm}
\begin{block}{Entrada e saída do modelo}
Dado um par \textbf{peptídeo--membrana}, predizer a probabilidade ou intensidade de interação funcional.
\end{block}

\vspace{-0.05cm}
\begin{center}
\resizebox{0.84\textwidth}{!}{%
\begin{tikzpicture}[node distance=0.65cm]
  \node[box,text width=2.3cm] (pep) {Peptídeo\ $p_i$};
  \node[box,text width=2.5cm,right=of pep] (mem) {Membrana\ $m_j$};
  \node[orangebox,text width=2.4cm,right=of mem] (model) {Modelo\ de IA};
  \node[bluebox,text width=2.7cm,right=of model] (out) {Score de\ interação};
  \draw[arrow] (pep) -- (mem);
  \draw[arrow] (mem) -- (model);
  \draw[arrow] (model) -- (out);
\end{tikzpicture}}
\end{center}

\vspace{-0.20cm}
\begin{equation*}
  f(p_i,m_j) \rightarrow y_{ij}
\end{equation*}
\vspace{-0.25cm}
\begin{itemize}
  \item $p_i$: representação do peptídeo.
  \item $m_j$: representação da membrana-alvo.
  \item $y_{ij}$: atividade, permeabilização, toxicidade ou seletividade.
\end{itemize}
\end{frame}

% -------------------------------------------------
\begin{frame}{Visão geral do pipeline}
\begin{center}
\resizebox{0.98\textwidth}{!}{%
\begin{tikzpicture}[node distance=0.95cm]
  \node[box] (data) {Coleta\ de dados};
  \node[box,right=of data] (cur) {Curadoria\ e padronização};
  \node[box,right=of cur] (rep) {Representação\ peptídeo/membrana};
  \node[orangebox,right=of rep] (model) {Modelagem\ de IA};
  \node[bluebox,right=of model] (rank) {Ranking\ de peptídeos};
  \node[box,right=of rank] (lab) {Validação\ in vitro};
  \node[box,below=1.1cm of model] (feed) {Realimentação\ da base};
  \draw[arrow] (data) -- (cur);
  \draw[arrow] (cur) -- (rep);
  \draw[arrow] (rep) -- (model);
  \draw[arrow] (model) -- (rank);
  \draw[arrow] (rank) -- (lab);
  \draw[dashedarrow] (lab) |- (feed);
  \draw[dashedarrow] (feed) -| (cur);
\end{tikzpicture}}
\end{center}

\vspace{0.3cm}
\begin{block}{Ideia principal}
A IA não substitui a bancada. Ela prioriza os candidatos mais promissores para reduzir custo, tempo e número de ensaios.
\end{block}
\end{frame}

% -------------------------------------------------
\begin{frame}{Etapa 1 - Construção da base de dados}
\begin{columns}[T]
\column{0.50\textwidth}
\textbf{Fontes de peptídeos}
\begin{itemize}
  \item Peptídeos escorpiônicos do projeto.
  \item Peptídeos antimicrobianos de bancos públicos.
  \item Peptídeos penetrantes celulares.
  \item Artigos com MIC, IC50, CC50, EC50 ou atividade antiviral.
\end{itemize}

\column{0.50\textwidth}
\textbf{Fontes de membranas e alvos}
\begin{itemize}
  \item Bactérias Gram+ e Gram-.
  \item Fungos e leveduras.
  \item Parasitas.
  \item Vírus envelopados.
  \item Células normais e tumorais.
\end{itemize}
\end{columns}

\vspace{0.3cm}
\begin{block}{Unidade da base}
Cada linha deve representar um par: \textbf{peptídeo + alvo + desfecho experimental}.
\end{block}
\end{frame}

% -------------------------------------------------
\begin{frame}{Endpoints experimentais: o que significam?}
\scriptsize
\setlength{\tabcolsep}{3pt}
\renewcommand{\arraystretch}{1.12}

\begin{columns}[T,onlytextwidth]
\column{0.49\textwidth}
\begin{tabular}{>{\raggedright\arraybackslash}p{1.15cm}>{\raggedright\arraybackslash}p{1.95cm}>{\raggedright\arraybackslash}p{3.45cm}}
\toprule
\textbf{End.} & \textbf{Significado} & \textbf{Leitura no projeto} \\
\midrule
MIC & Conc. Inibitória Mínima & Menor concentração que inibe crescimento microbiano visível. Menor MIC indica maior atividade. \\
\shortstack[l]{MBC/\\CBM} & Conc. Bactericida Mínima & Menor concentração capaz de matar bactérias. Diferencia ação bactericida de bacteriostática. \\
\shortstack[l]{MFC/\\CFM} & Conc. Fungicida Mínima & Menor concentração capaz de matar fungos ou leveduras. \\
\bottomrule
\end{tabular}

\column{0.49\textwidth}
\begin{tabular}{>{\raggedright\arraybackslash}p{0.95cm}>{\raggedright\arraybackslash}p{2.10cm}>{\raggedright\arraybackslash}p{3.50cm}}
\toprule
\textbf{End.} & \textbf{Significado} & \textbf{Leitura no projeto} \\
\midrule
IC50 & Conc. Inibitória 50\% & Concentração que reduz proliferação ou atividade celular em 50\%. Útil para células tumorais e parasitas. \\
CC50 & Conc. Citotóxica 50\% & Concentração que reduz viabilidade de células normais em 50\%. Maior CC50 indica menor toxicidade. \\
EC50 & Conc. Efetiva 50\% & Concentração que produz 50\% do efeito máximo, comum em ensaios antivirais. \\
SI & Índice de Seletividade & Razão entre toxicidade e atividade, por exemplo $SI=CC50/IC50$ ou $CC50/MIC$. Maior SI é melhor. \\
\bottomrule
\end{tabular}
\end{columns}

\vspace{0.15cm}
\begin{block}{Ideia central para a IA}
Os endpoints são transformados em rótulos ou valores-alvo para treinar o modelo de interação peptídeo--membrana.
\end{block}
\end{frame}

% -------------------------------------------------
\begin{frame}{Estrutura recomendada do dataset}
\scriptsize
\begin{center}
\setlength{\tabcolsep}{3pt}
\begin{tabular}{p{1.35cm}p{3.5cm}p{2.6cm}p{2.8cm}p{1.9cm}p{1.5cm}}
\toprule
\textbf{peptide\_id} & \textbf{sequence} & \textbf{target} & \textbf{target\_type} & \textbf{endpoint} & \textbf{value} \\
\midrule
P01 & \ttfamily FFSLIPSLVGGLISAFK & \textit{S. aureus} & Gram+ & MIC & 8 $\mu$M \\
P01 & \ttfamily FFSLIPSLVGGLISAFK & célula normal & mamífero normal & CC50 & 100 $\mu$M \\
P02 & \ttfamily FLGMIPGLIGGLISAFK & \textit{Candida} spp. & fungo & MIC & 16 $\mu$M \\
P03 & \ttfamily ... & Zika & vírus envelopado & redução viral & 70\% \\
\bottomrule
\end{tabular}
\end{center}

\vspace{0.10cm}
\begin{itemize}
  \item Cada linha representa um par \textbf{peptídeo + membrana-alvo + desfecho}.
  \item O mesmo peptídeo pode aparecer em várias linhas, uma para cada membrana-alvo.
  \item Isso permite aprender \alert{seletividade}: ativo contra alvo patológico e pouco tóxico contra célula normal.
\end{itemize}
\end{frame}

% -------------------------------------------------
\begin{frame}{Etapa 2 - Curadoria dos dados}
\begin{columns}[T]
\column{0.52\textwidth}
\begin{itemize}
  \item Padronizar sequências de aminoácidos.
  \item Remover duplicatas e registros inconsistentes.
  \item Converter unidades: $\mu$M, $\mu$g/mL, nM.
  \item Tratar modificações: amidado, acetilado, substituições.
  \item Padronizar nomes de organismos e linhagens.
  \item Definir classes: alta, média ou baixa atividade.
\end{itemize}

\column{0.48\textwidth}
\begin{block}{Exemplo de regra}
Para MIC, valores menores indicam maior atividade. Uma regra inicial pode ser:
\begin{itemize}
  \item MIC $\leq$ 10 $\mu$M: alta atividade.
  \item 10 $<$ MIC $\leq$ 50 $\mu$M: atividade moderada.
  \item MIC $>$ 50 $\mu$M: baixa atividade.
\end{itemize}
\end{block}
\end{columns}
\end{frame}

% -------------------------------------------------
\begin{frame}{Etapa 3 - Representação do peptídeo}
\begin{center}
\begin{tikzpicture}[node distance=1.0cm]
  \node[box] (seq) {Sequência\ de aminoácidos};
  \node[smallbox,below left=0.8cm and 0.15cm of seq] (desc) {Descritores\ físico-químicos};
  \node[smallbox,below right=0.8cm and 0.15cm of seq] (emb) {Embedding\ ProtBERT/ESM-2};
  \node[orangebox,below=2.1cm of seq] (vec) {Vetor final\ do peptídeo $p_i$};
  \draw[arrow] (seq) -- (desc);
  \draw[arrow] (seq) -- (emb);
  \draw[arrow] (desc) -- (vec);
  \draw[arrow] (emb) -- (vec);
\end{tikzpicture}
\end{center}

\vspace{0.15cm}
\begin{block}{Representação final}
\begin{equation*}
  p_i = [\text{descritores clássicos},\; \text{embedding da sequência}]
\end{equation*}
\end{block}
\end{frame}

% -------------------------------------------------
\begin{frame}{Descritores físico-químicos do peptídeo}
\small
\begin{center}
\begin{tabular}{p{3.5cm}p{8.0cm}}
\toprule
\textbf{Descritor} & \textbf{Interpretação biológica} \\
\midrule
Carga líquida & Interação eletrostática com membranas negativas. \\
Hidrofobicidade & Capacidade de inserção na bicamada lipídica. \\
Momento hidrofóbico & Organização anfipática, comum em peptídeos helicoidais. \\
Ponto isoelétrico & Comportamento de carga em função do pH. \\
Massa molecular & Tamanho e difusão do peptídeo. \\
Comprimento & Relacionado a estabilidade, poro e permeação. \\
Fração catiônica & Presença de Lys, Arg e His. \\
Fração hidrofóbica & Presença de Leu, Ile, Val, Phe, Trp etc. \\
\bottomrule
\end{tabular}
\end{center}
\end{frame}

% -------------------------------------------------
\begin{frame}{Embeddings de peptídeos}
\begin{block}{Ideia}
Modelos pré-treinados em grandes bases de proteínas transformam sequências em vetores densos que capturam padrões evolutivos, estruturais e funcionais.
\end{block}

\vspace{0.2cm}
\begin{center}
\begin{tikzpicture}[node distance=0.9cm]
  \node[box,text width=3.0cm] (seq) {FFSLIPSLVGGLISAFK};
  \node[orangebox,right=of seq,text width=2.8cm] (bert) {ProtBERT\ ou ESM-2};
  \node[bluebox,right=of bert,text width=3.2cm] (vec) {$[0.12,-0.04,\ldots,0.88]$};
  \draw[arrow] (seq) -- (bert);
  \draw[arrow] (bert) -- (vec);
\end{tikzpicture}
\end{center}

\vspace{0.2cm}
\begin{itemize}
  \item O embedding pode ser usado sozinho ou combinado com descritores clássicos.
  \item Para poucos dados, recomenda-se congelar o encoder e treinar apenas o classificador/regressor.
\end{itemize}
\end{frame}

% -------------------------------------------------
\begin{frame}{Etapa 4 - Representação da membrana-alvo}
\begin{block}{Ponto-chave}
A membrana não deve ser apenas uma classe como ``bactéria'' ou ``célula tumoral''. Ela deve ser descrita por propriedades físico-químicas e estruturais.
\end{block}

\vspace{0.2cm}
\begin{equation*}
  m_j = [T_j, C_j, A_j, L_j, S_j, R_j]
\end{equation*}

\begin{itemize}
  \item $T_j$: tipo de membrana.
  \item $C_j$: carga superficial.
  \item $A_j$: fração aniônica.
  \item $L_j$: composição lipídica.
  \item $S_j$: componentes estruturais.
  \item $R_j$: rigidez, fluidez e espessura.
\end{itemize}
\end{frame}

% -------------------------------------------------
\begin{frame}{Tipos de membrana-alvo}
\small
\begin{center}
\begin{tabular}{p{3.0cm}p{4.2cm}p{5.0cm}}
\toprule
\textbf{Classe} & \textbf{Exemplos} & \textbf{Uso no projeto} \\
\midrule
Gram-positiva & S. aureus, S. epidermidis & Atividade antibacteriana. \\
Gram-negativa & E. coli, P. aeruginosa & Atividade antibacteriana e barreira externa. \\
Fungo/levedura & Candida spp. & Atividade antifúngica. \\
Parasita & Trypanosoma cruzi & Atividade antiparasitária. \\
Vírus envelopado & Zika, HSV & Desestabilização de envelope viral. \\
Célula normal & linhagens não neoplásicas & Toxicidade indesejada. \\
Célula tumoral & linhagens cancerígenas & Atividade antiproliferativa. \\
\bottomrule
\end{tabular}
\end{center}
\end{frame}

% -------------------------------------------------
\begin{frame}{Descritores da membrana-alvo}
\small
\begin{center}
\begin{tabular}{p{3.7cm}p{8.2cm}}
\toprule
\textbf{Descritor} & \textbf{Importância} \\
\midrule
Carga superficial & Atrai ou repele peptídeos catiônicos. \\
Fração de lipídios aniônicos & Relacionada à seletividade para bactérias e células tumorais. \\
LPS & Característico de Gram-negativas; camada externa e carga negativa. \\
Peptidoglicano & Barreira estrutural em bactérias. \\
Ácido teicoico & Relevante em Gram-positivas. \\
Colesterol & Estabiliza membranas de mamíferos e pode reduzir permeabilização. \\
Ergosterol & Marcador importante de membranas fúngicas. \\
Envelope viral & Indica alvo lipídico em vírus envelopados. \\
Fluidez/rigidez & Controla deformação, inserção e formação de poros. \\
Espessura da bicamada & Afeta profundidade de inserção do peptídeo. \\
\bottomrule
\end{tabular}
\end{center}
\end{frame}

% -------------------------------------------------
\begin{frame}{Exemplo de codificação da membrana}
\small
\begin{center}
\begin{tabular}{p{2.4cm}rrrrrrr}
\toprule
\textbf{Alvo} & \textbf{Carga} & \textbf{Aniônica} & \textbf{LPS} & \textbf{PGN} & \textbf{Col.} & \textbf{Erg.} & \textbf{Env.} \\
\midrule
S. aureus & -0.80 & 0.60 & 0 & 1 & 0 & 0 & 0 \\
E. coli & -0.90 & 0.65 & 1 & 1 & 0 & 0 & 0 \\
Candida spp. & -0.40 & 0.25 & 0 & 0 & 0 & 1 & 0 \\
Zika & -0.30 & 0.20 & 0 & 0 & 1 & 0 & 1 \\
Célula normal & -0.10 & 0.10 & 0 & 0 & 1 & 0 & 0 \\
Célula tumoral & -0.40 & 0.25 & 0 & 0 & 1 & 0 & 0 \\
\bottomrule
\end{tabular}
\end{center}

\vspace{0.15cm}
\begin{itemize}
  \item Os valores iniciais podem ser aproximados da literatura.
  \item Depois, podem ser refinados com dados experimentais, lipossomos ou dinâmica molecular.
\end{itemize}
\end{frame}

% -------------------------------------------------
\begin{frame}{Níveis de representação da membrana}
\begin{columns}[T]
\column{0.33\textwidth}
\begin{block}{Nível 1\ simples}
\small
Tipo, carga, LPS, colesterol, ergosterol e fração aniônica.
\end{block}

\column{0.33\textwidth}
\begin{block}{Nível 2\ intermediário}
\small
Composição lipídica aproximada, fluidez, rigidez e espessura.
\end{block}

\column{0.33\textwidth}
\begin{block}{Nível 3\ avançado}
\small
Descritores de dinâmica molecular, lipossomos, potencial zeta e energia de inserção.
\end{block}
\end{columns}

\vspace{0.5cm}
\begin{center}
\begin{tikzpicture}[node distance=1.0cm]
  \node[smallbox] (n1) {rápido\ e robusto};
  \node[smallbox,right=of n1] (n2) {melhor\ biologia};
  \node[smallbox,right=of n2] (n3) {maior\ mecanismo};
  \draw[arrow] (n1) -- (n2);
  \draw[arrow] (n2) -- (n3);
\end{tikzpicture}
\end{center}
\end{frame}

% -------------------------------------------------
\begin{frame}{Etapa 5 - Construção dos pares peptídeo--membrana}
\begin{block}{Entrada final do modelo}
\begin{equation*}
  x_{ij} = [p_i, m_j, g(p_i,m_j)]
\end{equation*}
\end{block}

\begin{itemize}
  \item $p_i$: descritores e embedding do peptídeo.
  \item $m_j$: descritores da membrana.
  \item $g(p_i,m_j)$: atributos cruzados de interação.
\end{itemize}

\vspace{0.2cm}
\begin{center}
\begin{tikzpicture}[node distance=0.9cm]
  \node[box] (p) {Vetor\ peptídeo};
  \node[box,right=of p] (m) {Vetor\ membrana};
  \node[orangebox,right=of m] (g) {Atributos\ cruzados};
  \node[bluebox,right=of g] (x) {Entrada\ $x_{ij}$};
  \draw[arrow] (p) -- (m);
  \draw[arrow] (m) -- (g);
  \draw[arrow] (g) -- (x);
\end{tikzpicture}
\end{center}
\end{frame}

% -------------------------------------------------
\begin{frame}{Atributos cruzados peptídeo--membrana}
\begin{columns}[T]
\column{0.54\textwidth}
\begin{itemize}
  \item Compatibilidade eletrostática.
  \item Compatibilidade hidrofóbica.
  \item Penalização por colesterol.
  \item Afinidade por membranas aniônicas.
  \item Potencial de seletividade.
\end{itemize}

\column{0.46\textwidth}
\begin{block}{Ideia}
\small
Além de representar peptídeo e membrana separadamente, podemos criar descritores que expressem \textbf{como os dois combinam entre si}.
\end{block}
\end{columns}

\vspace{0.25cm}
\begin{center}
\begin{tikzpicture}[node distance=0.85cm]
  \node[box] (p) {Features\ do peptídeo};
  \node[box,right=of p] (m) {Features\ da membrana};
  \node[orangebox,right=of m] (g) {Features\ de interação};
  \draw[arrow] (p) -- (m);
  \draw[arrow] (m) -- (g);
\end{tikzpicture}
\end{center}
\end{frame}

% -------------------------------------------------
\begin{frame}{PMI: Peptide--Membrane Interaction Index}
\begin{block}{Objetivo do PMI}
Construir um índice simples e interpretável para resumir a compatibilidade entre um peptídeo e uma membrana-alvo antes ou junto do modelo de IA.
\end{block}

\vspace{0.2cm}
\begin{equation*}
PMI = \alpha Q_p|Q_m| + \beta H_pH_m + \gamma \mu H_p - \delta Chol_m
\end{equation*}

\vspace{0.1cm}
\begin{itemize}
  \item Valores maiores sugerem maior potencial de interação, inserção ou perturbação da membrana.
  \item O PMI pode ser usado como feature adicional ou como baseline interpretável.
  \item Os pesos $\alpha$, $\beta$, $\gamma$ e $\delta$ podem ser definidos empiricamente ou aprendidos pelo modelo.
\end{itemize}
\end{frame}

% -------------------------------------------------
\begin{frame}{PMI: interpretação dos termos}
\small
\begin{center}
\begin{tabular}{p{2.4cm}p{4.0cm}p{5.7cm}}
\toprule
\textbf{Termo} & \textbf{Representa} & \textbf{Intuição biológica} \\
\midrule
$Q_p|Q_m|$ & compatibilidade eletrostática & Peptídeos catiônicos tendem a interagir com membranas mais negativas. \\
$H_pH_m$ & compatibilidade hidrofóbica & Peptídeos hidrofóbicos podem se inserir melhor em ambientes lipídicos. \\
$\mu H_p$ & momento hidrofóbico & Captura anfipaticidade, importante para hélices que interagem com bicamadas. \\
$Chol_m$ & fração de colesterol & Pode estabilizar membranas de mamíferos e reduzir permeabilização por alguns peptídeos. \\
\bottomrule
\end{tabular}
\end{center}

\vspace{0.15cm}
\begin{block}{Leitura prática}
Um peptídeo bom deve ter PMI alto para a membrana patológica e PMI baixo para a célula normal.
\end{block}
\end{frame}

% -------------------------------------------------
\begin{frame}{PMI e seletividade}
\begin{block}{Problema biológico}
Não basta interagir fortemente com membranas. É necessário interagir preferencialmente com o alvo patológico e pouco com células normais.
\end{block}

\vspace{0.2cm}
\begin{equation*}
PMI_{sel} = PMI_{alvo\ patológico} - PMI_{célula\ normal}
\end{equation*}

\vspace{0.1cm}
\begin{center}
\resizebox{0.88\textwidth}{!}{%
\begin{tikzpicture}[node distance=0.85cm]
  \node[box] (pep) {Peptídeo};
  \node[orangebox,right=of pep] (target) {$PMI$\ alvo patológico};
  \node[box,below=0.8cm of target] (normal) {$PMI$\ célula normal};
  \node[bluebox,right=1.1cm of target] (sel) {$PMI_{sel}$\ seletividade};
  \draw[arrow] (pep) -- (target);
  \draw[arrow] (pep) |- (normal);
  \draw[arrow] (target) -- (sel);
  \draw[arrow] (normal) -| (sel);
\end{tikzpicture}}
\end{center}

\vspace{0.1cm}
\begin{itemize}
  \item $PMI_{sel} > 0$: favorece alvo patológico.
  \item $PMI_{sel} \approx 0$: baixa seletividade prevista.
  \item $PMI_{sel} < 0$: risco de maior interação com célula normal.
\end{itemize}
\end{frame}

% -------------------------------------------------
\begin{frame}{PMI: como usar no pipeline}
\small
\begin{columns}[T]
\column{0.50\textwidth}
\textbf{Uso como baseline}
\begin{itemize}
  \item Calcular o PMI para todos os pares peptídeo--membrana.
  \item Ordenar os candidatos por $PMI_{sel}$.
  \item Comparar com resultados experimentais.
\end{itemize}

\column{0.50\textwidth}
\textbf{Uso como feature}
\begin{itemize}
  \item Adicionar PMI ao vetor $x_{ij}$.
  \item Permitir que RF, XGBoost ou MLP aprendam pesos melhores.
  \item Usar SHAP para verificar a importância do PMI.
\end{itemize}
\end{columns}

\vspace{0.25cm}
\begin{block}{Cuidado}
O PMI é uma simplificação. Ele não substitui embeddings, composição lipídica detalhada, dinâmica molecular ou validação experimental.
\end{block}
\end{frame}

% -------------------------------------------------
\begin{frame}{Etapa 6 - Modelos de IA}
\begin{center}
\begin{tikzpicture}[node distance=1.0cm]
  \node[box,text width=3.0cm] (l1) {Nível 1\ Baselines clássicos};
  \node[orangebox,text width=3.2cm,right=of l1] (l2) {Nível 2\ Embeddings + MLP};
  \node[bluebox,text width=3.4cm,right=of l2] (l3) {Nível 3\ Modelo multimodal};
  \draw[arrow] (l1) -- (l2);
  \draw[arrow] (l2) -- (l3);
\end{tikzpicture}
\end{center}

\vspace{0.3cm}
\begin{itemize}
  \item \textbf{Nível 1}: Random Forest, XGBoost, SVM e regressão logística.
  \item \textbf{Nível 2}: embeddings ProtBERT/ESM-2 combinados com descritores.
  \item \textbf{Nível 3}: encoder de peptídeo + encoder de membrana + fusão multimodal.
\end{itemize}
\end{frame}

% -------------------------------------------------
\begin{frame}{Por que o modelo deve ser multimodal?}
\begin{block}{Motivação}
O problema peptídeo--membrana combina diferentes tipos de informação: sequência, descritores físico-químicos, propriedades da membrana e endpoints experimentais.
\end{block}

\vspace{0.25cm}
\begin{center}
\resizebox{0.92\textwidth}{!}{%
\begin{tikzpicture}[node distance=0.85cm]
  \node[box] (seq) {Sequência\ do peptídeo};
  \node[box,right=of seq] (desc) {Descritores\ do peptídeo};
  \node[box,right=of desc] (mem) {Descritores\ da membrana};
  \node[orangebox,right=of mem] (fusion) {Fusão\ multimodal};
  \node[bluebox,right=of fusion] (pred) {Predições};
  \draw[arrow] (seq) -- (desc);
  \draw[arrow] (desc) -- (mem);
  \draw[arrow] (mem) -- (fusion);
  \draw[arrow] (fusion) -- (pred);
\end{tikzpicture}}
\end{center}

\vspace{0.15cm}
\begin{itemize}
  \item A sequência captura padrões de resíduos e motivos funcionais.
  \item Os descritores clássicos capturam propriedades conhecidas e interpretáveis.
  \item A membrana define o contexto biológico da interação.
\end{itemize}
\end{frame}

% -------------------------------------------------
\begin{frame}{Arquitetura multimodal proposta}
\begin{center}
\resizebox{0.94\textwidth}{!}{%
\begin{tikzpicture}[node distance=0.75cm]
  \node[box] (seq) {Sequência\ peptídica};
  \node[box,below=of seq] (desc) {Descritores\ do peptídeo};
  \node[box,below=of desc] (mem) {Descritores\ da membrana};

  \node[orangebox,right=1.1cm of seq] (enc1) {Encoder\ ProtBERT/ESM-2};
  \node[orangebox,right=1.1cm of desc] (enc2) {MLP\ peptídeo};
  \node[orangebox,right=1.1cm of mem] (enc3) {MLP\ membrana};

  \node[bluebox,right=1.2cm of enc2] (fusion) {Fusão\ multimodal};
  \node[orangebox,right=1.0cm of fusion] (head) {Cabeças\ preditivas};
  \node[bluebox,right=1.0cm of head] (out) {Atividade\ Toxicidade\ Seletividade};

  \draw[arrow] (seq) -- (enc1);
  \draw[arrow] (desc) -- (enc2);
  \draw[arrow] (mem) -- (enc3);
  \draw[arrow] (enc1) -- (fusion);
  \draw[arrow] (enc2) -- (fusion);
  \draw[arrow] (enc3) -- (fusion);
  \draw[arrow] (fusion) -- (head);
  \draw[arrow] (head) -- (out);
\end{tikzpicture}}
\end{center}
\end{frame}

% -------------------------------------------------
\begin{frame}{Fusão multimodal: opções práticas}
\small
\begin{center}
\begin{tabular}{p{3.2cm}p{4.1cm}p{4.8cm}}
\toprule
\textbf{Estratégia} & \textbf{Como funciona} & \textbf{Quando usar} \\
\midrule
Concatenação & Junta os vetores em uma única entrada & Primeira versão, simples e robusta. \\
Fusão com MLP & Cada modalidade passa por uma MLP antes da fusão & Quando descritores têm escalas diferentes. \\
Atenção cruzada & O vetor da membrana modula quais partes do peptídeo importam & Quando houver dados suficientes. \\
Fusão multitarefa & Um tronco comum alimenta várias saídas & Para prever atividade, toxicidade e seletividade juntas. \\
\bottomrule
\end{tabular}
\end{center}

\vspace{0.15cm}
\begin{block}{Recomendação inicial}
Começar com concatenação + MLP, usando embeddings congelados, e evoluir para atenção cruzada quando a base crescer.
\end{block}
\end{frame}

% -------------------------------------------------
\begin{frame}{Cabeças preditivas do modelo multimodal}
\begin{center}
\resizebox{0.82\textwidth}{!}{%
\begin{tikzpicture}[node distance=0.9cm]
  \node[orangebox,text width=3.2cm] (z) {Representação\ fusionada $z_{ij}$};
  \node[box,right=1.1cm of z,yshift=1.0cm] (h1) {Cabeça 1\ atividade};
  \node[box,right=1.1cm of z] (h2) {Cabeça 2\ toxicidade};
  \node[box,right=1.1cm of z,yshift=-1.0cm] (h3) {Cabeça 3\ seletividade};
  \node[bluebox,right=1.0cm of h1] (o1) {MIC, IC50\ ou classe};
  \node[bluebox,right=1.0cm of h2] (o2) {CC50\ ou toxicidade};
  \node[bluebox,right=1.0cm of h3] (o3) {$SI$ ou\ ranking};
  \draw[arrow] (z) -- (h1);
  \draw[arrow] (z) -- (h2);
  \draw[arrow] (z) -- (h3);
  \draw[arrow] (h1) -- (o1);
  \draw[arrow] (h2) -- (o2);
  \draw[arrow] (h3) -- (o3);
\end{tikzpicture}}
\end{center}

\vspace{0.1cm}
\begin{itemize}
  \item A saída multitarefa força o modelo a aprender atividade e segurança simultaneamente.
  \item O ranking final pode combinar as três cabeças preditivas.
\end{itemize}
\end{frame}

% -------------------------------------------------
\begin{frame}{Tipos de saída do modelo}
\small
\begin{center}
\begin{tabular}{p{3.0cm}p{4.5cm}p{4.5cm}}
\toprule
\textbf{Tarefa} & \textbf{Saída} & \textbf{Exemplo} \\
\midrule
Classificação binária & Interage ou não interage & membrana-ativo / não ativo \\
Classificação ordinal & Baixa, média ou alta interação & classe de atividade \\
Regressão & Valor contínuo & MIC, IC50, CC50, EC50 \\
Ranqueamento & Ordem dos candidatos & Top-5 peptídeos para bancada \\
Multitarefa & Várias saídas simultâneas & atividade, toxicidade, seletividade \\
\bottomrule
\end{tabular}
\end{center}

\vspace{0.2cm}
\begin{block}{Recomendação}
Começar com classificação/ranking e evoluir para regressão e multitarefa quando houver mais dados.
\end{block}
\end{frame}

% -------------------------------------------------
\begin{frame}{Etapa 7 - Treinamento e validação}
\begin{columns}[T]
\column{0.50\textwidth}
\textbf{Divisão dos dados}
\begin{itemize}
  \item Treino.
  \item Validação.
  \item Teste interno.
  \item Teste externo com literatura.
  \item Teste prospectivo com bancada.
\end{itemize}

\column{0.50\textwidth}
\textbf{Métricas}
\begin{itemize}
  \item Classificação: acurácia, precisão, recall, F1, ROC-AUC.
  \item Regressão: MAE, RMSE, $R^2$.
  \item Ranking: top-k accuracy, enrichment factor.
  \item Seletividade: correlação com índice experimental.
\end{itemize}
\end{columns}
\end{frame}

% -------------------------------------------------
\begin{frame}{Cuidado metodológico: vazamento de informação}
\begin{block}{Problema comum}
Se peptídeos muito parecidos aparecem simultaneamente no treino e no teste, o desempenho pode ficar artificialmente alto.
\end{block}

\vspace{0.2cm}
\begin{itemize}
  \item Fazer split por cluster de similaridade, não apenas split aleatório.
  \item Separar peptídeos análogos próximos.
  \item Usar teste externo com dados não usados na curadoria.
  \item Reportar desempenho por tipo de alvo: Gram+, Gram-, fungo, célula normal etc.
\end{itemize}
\end{frame}

% -------------------------------------------------
\begin{frame}{Etapa 8 - Ranqueamento dos peptídeos}
\begin{block}{Critério recomendado}
O melhor peptídeo não é necessariamente o mais ativo. É o que combina atividade alta contra o alvo e baixa toxicidade em célula normal.
\end{block}

\vspace{0.2cm}
\begin{equation*}
  Score_{final} = Score_{alvo} - \lambda \cdot Score_{tox} + \eta \cdot Score_{seletividade}
\end{equation*}

\begin{itemize}
  \item $Score_{alvo}$: interação prevista com bactéria, fungo, parasita, vírus ou célula tumoral.
  \item $Score_{tox}$: interação prevista com célula normal.
  \item $Score_{seletividade}$: diferença entre alvo patológico e célula normal.
\end{itemize}
\end{frame}

% -------------------------------------------------
\begin{frame}{Exemplo de ranking para bancada}
\small
\begin{center}
\begin{tabular}{p{2.0cm}rrrrp{2.6cm}}
\toprule
\textbf{Peptídeo} & \textbf{Alvo} & \textbf{Tox.} & \textbf{Sel.} & \textbf{Score} & \textbf{Decisão} \\
\midrule
P04 & 0.91 & 0.18 & 0.73 & 0.88 & Testar primeiro \\
P07 & 0.86 & 0.22 & 0.64 & 0.80 & Testar \\
P02 & 0.78 & 0.55 & 0.23 & 0.52 & Cuidado \\
P09 & 0.65 & 0.12 & 0.53 & 0.70 & Alternativa \\
P03 & 0.93 & 0.89 & 0.04 & 0.30 & Evitar por toxicidade \\
\bottomrule
\end{tabular}
\end{center}

\vspace{0.2cm}
\begin{itemize}
  \item Este ranking é um exemplo didático.
  \item Os valores reais viriam do modelo treinado e da validação experimental.
\end{itemize}
\end{frame}

% -------------------------------------------------
\begin{frame}{Etapa 9 - Validação experimental}
\begin{center}
\resizebox{0.96\textwidth}{!}{%
\begin{tikzpicture}[node distance=0.8cm]
  \node[box] (rank) {Peptídeos\ priorizados};
  \node[box,right=of rank] (anti) {MIC/MBC\ bactérias};
  \node[box,right=of anti] (fung) {MIC/MFC\ fungos};
  \node[box,right=of fung] (para) {T. cruzi\ inibição};
  \node[box,right=of para] (viral) {Zika/HSV\ redução viral};
  \node[box,right=of viral] (tox) {CC50/IC50\ células};
  \node[orangebox,below=1.0cm of para] (mev) {MEV\ evidência morfológica};
  \draw[arrow] (rank) -- (anti);
  \draw[arrow] (anti) -- (fung);
  \draw[arrow] (fung) -- (para);
  \draw[arrow] (para) -- (viral);
  \draw[arrow] (viral) -- (tox);
  \draw[dashedarrow] (anti) |- (mev);
  \draw[dashedarrow] (fung) |- (mev);
  \draw[dashedarrow] (para) -- (mev);
  \draw[dashedarrow] (viral) |- (mev);
\end{tikzpicture}}
\end{center}

\vspace{0.2cm}
\begin{block}{Resultado esperado}
Confirmar se o modelo preditivo realmente seleciona peptídeos com atividade e seletividade.
\end{block}
\end{frame}

% -------------------------------------------------
\begin{frame}{Etapa 10 - Ciclo de active learning}
\begin{center}
\begin{tikzpicture}[node distance=1.0cm]
  \node[box] (base) {Base\ inicial};
  \node[box,right=of base] (model) {Modelo\ preditivo};
  \node[orangebox,right=of model] (rank) {Ranking\ dos candidatos};
  \node[box,below=1.0cm of rank] (lab) {Ensaios\ in vitro};
  \node[bluebox,left=of lab] (new) {Novos\ dados};
  \node[box,left=of new] (ret) {Re-treinamento};
  \draw[arrow] (base) -- (model);
  \draw[arrow] (model) -- (rank);
  \draw[arrow] (rank) -- (lab);
  \draw[arrow] (lab) -- (new);
  \draw[arrow] (new) -- (ret);
  \draw[arrow] (ret) -- (base);
\end{tikzpicture}
\end{center}

\vspace{0.2cm}
\begin{block}{Aprendizado ativo}
A cada rodada experimental, o modelo aprende melhor quais propriedades aumentam atividade, seletividade e segurança.
\end{block}
\end{frame}

% -------------------------------------------------
\begin{frame}{Explicabilidade do modelo}
\begin{columns}[T]
\column{0.50\textwidth}
\textbf{Ferramentas}
\begin{itemize}
  \item SHAP.
  \item Permutation importance.
  \item Atenção em Transformers.
  \item Análise de resíduos importantes.
\end{itemize}

\column{0.50\textwidth}
\textbf{Perguntas biológicas}
\begin{itemize}
  \item A carga positiva aumenta interação?
  \item Hidrofobicidade excessiva aumenta toxicidade?
  \item Quais resíduos favorecem seletividade?
  \item O modelo distingue célula normal de tumor?
\end{itemize}
\end{columns}

\vspace{0.3cm}
\begin{block}{Objetivo da XAI}
Transformar predição em conhecimento de relação estrutura--atividade.
\end{block}
\end{frame}

% -------------------------------------------------
\begin{frame}{Entregáveis computacionais}
\begin{itemize}
  \item Base de dados curada de peptídeos, membranas e desfechos experimentais.
  \item Scripts de extração de descritores físico-químicos.
  \item Geração de embeddings por ProtBERT, ESM-2 ou ProtT5.
  \item Tabela de representação de membranas-alvo.
  \item Modelos baseline e multimodais treinados.
  \item Ranking dos peptídeos do projeto.
  \item Relatório com métricas, análise de erro e explicabilidade.
  \item Pipeline reprodutível em Python, preferencialmente com notebooks e scripts versionados.
\end{itemize}
\end{frame}

% -------------------------------------------------
\begin{frame}{Plano de execução sugerido}
\small
\begin{center}
\begin{tabular}{p{1.2cm}p{5.5cm}p{5.5cm}}
\toprule
\textbf{Fase} & \textbf{Atividade} & \textbf{Produto} \\
\midrule
1 & Revisão de bases e coleta inicial & planilha padronizada \\
2 & Cálculo de descritores dos peptídeos & tabela de features \\
3 & Construção da tabela de membranas & vetor $m_j$ para cada alvo \\
4 & Geração de embeddings & matriz de embeddings \\
5 & Baseline RF/XGBoost/SVM & métricas iniciais \\
6 & Modelo multimodal & predição peptídeo--membrana \\
7 & Ranking dos candidatos & lista priorizada para bancada \\
8 & XAI e relatório & interpretação biológica \\
\bottomrule
\end{tabular}
\end{center}
\end{frame}

% -------------------------------------------------
\begin{frame}{Stack computacional sugerido}
\begin{columns}[T]
\column{0.50\textwidth}
\textbf{Python e dados}
\begin{itemize}
  \item pandas, numpy, scipy.
  \item scikit-learn.
  \item xgboost.
  \item pyarrow/parquet.
\end{itemize}

\column{0.50\textwidth}
\textbf{Deep learning e bioinformática}
\begin{itemize}
  \item PyTorch.
  \item Hugging Face Transformers.
  \item ESM / ProtBERT / ProtT5.
  \item Biopython.
  \item SHAP.
\end{itemize}
\end{columns}

\vspace{0.3cm}
\begin{block}{Organização}
Separar o projeto em módulos: dados, features, embeddings, modelos, avaliação, explicabilidade e ranking.
\end{block}
\end{frame}

% -------------------------------------------------
% -------------------------------------------------
\begin{frame}{Possíveis entregáveis do projeto}
\begin{center}
\begin{tikzpicture}[node distance=0.9cm]
  \node[box, text width=3.2cm] (artigos) {\textbf{Artigos}\\modelagem, validação e aplicação};
  \node[box, text width=3.2cm, right=1.1cm of artigos] (datasets) {\textbf{Datasets}\\dados curados e reprodutíveis};
  \node[box, text width=3.2cm, right=1.1cm of datasets] (plataforma) {\textbf{Plataforma Web}\\predição, ranking e visualização};
  \node[smallbox, text width=3.2cm, below=0.8cm of artigos] (a1) {metodologia PepMem-AI};
  \node[smallbox, text width=3.2cm, below=0.8cm of datasets] (d1) {peptídeos + membranas + endpoints};
  \node[smallbox, text width=3.2cm, below=0.8cm of plataforma] (p1) {interface para alunos e pesquisadores};
  \draw[arrow] (artigos) -- (a1);
  \draw[arrow] (datasets) -- (d1);
  \draw[arrow] (plataforma) -- (p1);
\end{tikzpicture}
\end{center}
\vspace{0.2cm}
\begin{block}{Visão estratégica}
Cada entregável deve reforçar os demais: os dados alimentam os modelos, os modelos geram artigos, e a plataforma torna o pipeline utilizável e demonstrável.
\end{block}
\end{frame}

% -------------------------------------------------
\begin{frame}{Sugestões de artigos: linha 1}
\begin{block}{Artigo metodológico}
\textbf{PepMem-AI: um pipeline multimodal para predição de interação peptídeo--membrana}
\end{block}
\begin{columns}[T]
\column{0.48\textwidth}
\textbf{Ideia central}
\begin{itemize}
  \item Formalizar o problema como pares peptídeo--membrana.
  \item Comparar descritores clássicos, embeddings e modelos multimodais.
  \item Avaliar classificação, regressão e ranking.
\end{itemize}
\column{0.52\textwidth}
\textbf{Resultados esperados}
\begin{itemize}
  \item Baselines: RF, SVM, XGBoost e MLP.
  \item Modelos com ProtBERT/ESM-2.
  \item Métricas: ROC-AUC, F1, RMSE, R$^2$ e top-k.
  \item Análise XAI dos descritores relevantes.
\end{itemize}
\end{columns}
\end{frame}

% -------------------------------------------------
\begin{frame}{Sugestões de artigos: linhas 2 e 3}
\begin{columns}[T]
\column{0.50\textwidth}
\begin{block}{Artigo experimental-computacional}
\textbf{Validação in vitro de peptídeos priorizados por IA}
\end{block}
\begin{itemize}
  \item Usar o ranking do modelo para selecionar candidatos.
  \item Validar MIC, MBC/CBM, IC50, CC50 e atividade antiviral.
  \item Comparar predição versus bancada.
  \item Discutir seletividade contra células normais.
\end{itemize}
\column{0.50\textwidth}
\begin{block}{Artigo de interpretabilidade}
\textbf{Determinantes físico-químicos da interação peptídeo--membrana}
\end{block}
\begin{itemize}
  \item Usar SHAP e análise de importância.
  \item Explorar carga, hidrofobicidade e momento hidrofóbico.
  \item Relacionar PMI com endpoints experimentais.
  \item Gerar hipóteses para novos análogos.
\end{itemize}
\end{columns}
\end{frame}

% -------------------------------------------------
\begin{frame}{Datasets possíveis}
\begin{columns}[T]
\column{0.50\textwidth}
\textbf{Dataset 1: PepMem-Base}
\begin{itemize}
  \item Sequências de peptídeos escorpiônicos.
  \item Descritores físico-químicos.
  \item Embeddings ProtBERT/ESM-2.
  \item Origem: APD3, CAMP, UniProt e literatura.
\end{itemize}
\vspace{0.1cm}
\textbf{Dataset 2: Membrane-Targets}
\begin{itemize}
  \item Representações de membranas-alvo.
  \item Tipo, carga, fração aniônica, colesterol, ergosterol, LPS e peptidoglicano.
\end{itemize}
\column{0.50\textwidth}
\textbf{Dataset 3: PepMem-Endpoints}
\begin{itemize}
  \item Pares peptídeo--membrana.
  \item Endpoints: MIC, MBC/CBM, IC50, CC50, EC50 e SI.
  \item Dados da literatura e da bancada.
\end{itemize}
\vspace{0.2cm}
\begin{block}{Formato recomendado}
CSV/Parquet + dicionário de dados + scripts de reprodução + licença + DOI em Mendeley Data, Zenodo ou Figshare.
\end{block}
\end{columns}
\end{frame}

% -------------------------------------------------
\begin{frame}{Estrutura mínima do dataset}
\scriptsize
\begin{tabular}{p{0.20\textwidth} p{0.29\textwidth} p{0.43\textwidth}}
\toprule
\textbf{Grupo} & \textbf{Campos} & \textbf{Finalidade} \\
\midrule
Identificação & peptide\_id, sequence, target, target\_type & rastrear o par peptídeo--membrana \\
Descritores do peptídeo & net\_charge, hydrophobicity, hydrophobic\_moment, pI, length & representar propriedades físico-químicas \\
Embeddings & esm2\_vector, protbert\_vector & representar a sequência em espaço vetorial \\
Membrana & surface\_charge, anionic\_fraction, cholesterol, ergosterol, LPS, PG, PE, CL & representar a membrana-alvo \\
Endpoint & endpoint, value, unit, assay, reference & registrar o desfecho experimental \\
Controle de qualidade & source, duplicate\_flag, confidence, split & permitir curadoria, reprodutibilidade e treino/teste \\
\bottomrule
\end{tabular}
\vspace{0.3cm}
\begin{block}{Entregável de dados}
Uma versão inicial pode conter apenas dados de literatura. A versão final deve incorporar os resultados experimentais gerados no projeto.
\end{block}
\end{frame}

% -------------------------------------------------
\begin{frame}{Plataforma web: visão geral}
\begin{center}
\begin{tikzpicture}[node distance=0.65cm]
  \node[box, text width=2.4cm] (input) {Entrada\\sequência + alvo};
  \node[box, text width=2.6cm, right=0.8cm of input] (features) {Features\\descritores + embeddings};
  \node[box, text width=2.5cm, right=0.8cm of features] (model) {Modelo\\multimodal};
  \node[box, text width=2.5cm, right=0.8cm of model] (output) {Saída\\ranking + relatório};
  \draw[arrow] (input) -- (features);
  \draw[arrow] (features) -- (model);
  \draw[arrow] (model) -- (output);
  \node[smallbox, text width=2.7cm, below=0.8cm of features] (db) {Banco de dados\\peptídeos, membranas e endpoints};
  \node[smallbox, text width=2.7cm, below=0.8cm of model] (xai) {XAI\\SHAP, importância e PMI};
  \draw[dashedarrow] (db) -- (features);
  \draw[dashedarrow] (model) -- (xai);
  \draw[dashedarrow] (xai) -- (output);
\end{tikzpicture}
\end{center}
\vspace{0.2cm}
\begin{block}{Objetivo da plataforma}
Permitir que o usuário submeta uma sequência peptídica, escolha uma membrana-alvo e receba predições de atividade, toxicidade, seletividade e explicabilidade.
\end{block}
\end{frame}

% -------------------------------------------------
\begin{frame}{Plataforma web: módulos principais}
\begin{columns}[T]
\column{0.48\textwidth}
\textbf{Módulos do usuário}
\begin{itemize}
  \item Cadastro ou upload de peptídeos.
  \item Escolha da membrana-alvo.
  \item Visualização do ranking.
  \item Relatório com justificativa da predição.
\end{itemize}
\vspace{0.2cm}
\textbf{Módulos científicos}
\begin{itemize}
  \item Cálculo de descritores.
  \item Geração de embeddings.
  \item Predição multimodal.
  \item XAI e PMI.
\end{itemize}
\column{0.52\textwidth}
\textbf{Saídas da plataforma}
\begin{itemize}
  \item Score de interação peptídeo--membrana.
  \item Score de toxicidade prevista.
  \item Índice de seletividade previsto.
  \item Top-k peptídeos para validação in vitro.
  \item Download do relatório em PDF/CSV.
\end{itemize}
\begin{block}{Diferencial}
A plataforma transforma o pipeline em uma ferramenta prática para triagem, ensino, colaboração e geração de novos experimentos.
\end{block}
\end{columns}
\end{frame}

% -------------------------------------------------
\begin{frame}{Roadmap de entregáveis}
\begin{center}
\begin{tikzpicture}[node distance=0.65cm]
  \node[box, text width=2.1cm] (m1) {M1--M3\\base inicial};
  \node[box, text width=2.1cm, right=0.65cm of m1] (m2) {M4--M6\\baselines};
  \node[box, text width=2.1cm, right=0.65cm of m2] (m3) {M7--M9\\multimodal};
  \node[box, text width=2.1cm, right=0.65cm of m3] (m4) {M10--M12\\validação};
  \node[box, text width=2.1cm, right=0.65cm of m4] (m5) {Final\\artigos + web};
  \draw[arrow] (m1) -- (m2);
  \draw[arrow] (m2) -- (m3);
  \draw[arrow] (m3) -- (m4);
  \draw[arrow] (m4) -- (m5);
\end{tikzpicture}
\end{center}
\vspace{0.3cm}
\begin{itemize}
  \item \textbf{Curto prazo:} dataset curado, baselines e primeira versão do ranking.
  \item \textbf{Médio prazo:} modelo multimodal, PMI, XAI e validação com literatura.
  \item \textbf{Final:} validação experimental, artigos, dataset público e protótipo web.
\end{itemize}
\end{frame}

\begin{frame}{Resumo final}
\begin{enumerate}
  \item O problema central é \alert{predição de interação peptídeo--membrana}.
  \item A entrada é um par: peptídeo + membrana-alvo.
  \item O peptídeo deve ser representado por descritores clássicos e embeddings.
  \item A membrana deve ser representada por carga, composição, componentes estruturais e propriedades físicas.
  \item O modelo deve prever atividade, toxicidade e seletividade.
  \item A saída mais útil é um ranking para validação experimental.
  \item Os dados da bancada devem realimentar o modelo em um ciclo de active learning.
\end{enumerate}
\end{frame}

% -------------------------------------------------
\begin{frame}{Mensagem final}
\begin{block}{Ideia principal}
O objetivo não é apenas treinar um modelo para dizer se um peptídeo é ativo. O objetivo é aprender \textbf{por que}, \textbf{contra qual membrana} e \textbf{com que seletividade} ele pode funcionar.
\end{block}

\vspace{0.3cm}
\begin{center}
\Large\bfseries\color{InovGreen}
Peptídeo + Membrana + IA + Bancada = Bioprospecção orientada por dados
\end{center}
\end{frame}

% -------------------------------------------------
\begin{frame}{Referências de apoio}
\small
\begin{itemize}
  \item Lee et al. Mapping membrane activity in undiscovered peptide sequence space using machine learning. \textit{PNAS}, 2016.
  \item Damiati et al. Novel machine learning application for prediction of membrane insertion potential of cell-penetrating peptides. \textit{International Journal of Pharmaceutics}, 2019.
  \item de Oliveira et al. Predicting cell-penetrating peptides using machine learning algorithms and navigating in their chemical space. \textit{Scientific Reports}, 2021.
  \item Tran et al. Using molecular dynamics simulations to prioritize and understand AI-generated cell penetrating peptides. \textit{Scientific Reports}, 2021.
  \item Cao et al. Multi\_CycGT: A Deep Learning-Based Multimodal Model for Predicting the Membrane Permeability of Cyclic Peptides. \textit{Journal of Medicinal Chemistry}, 2024.
\end{itemize}
\end{frame}

% -------------------------------------------------
\begin{frame}[plain]
  \begin{tikzpicture}[remember picture,overlay]
    \node[anchor=east] at ([xshift=-0.40cm]current page.east) {\includegraphics[width=0.31\paperwidth]{/mnt/data/inovai_instagram_card_narrow.png}};
  \end{tikzpicture}
  \begin{center}
    \vspace{-0.20cm}
    {\Huge\bfseries\color{InovGreen} Obrigado!}\par
    \vspace{0.25cm}
    \includegraphics[width=0.14\paperwidth]{/mnt/data/tech_meets_biology_logo_design.png}\\[0.15cm]
    {\Large PepMem-AI}\par
    \vspace{0.12cm}
    {\large Pipeline de IA para interação peptídeo--membrana}\par
    \vspace{0.18cm}
    {\large InovAI Lab \; | \; IMD/UFRN}\par
  \end{center}
\end{frame}

\end{document}
